%% LyX 2.5.1 created this file.  For more info, see https://www.lyx.org/.
%% Do not edit unless you really know what you are doing.
\documentclass[twoside,onecolumn,american,compsoc]{IEEEtran}
\usepackage[T1]{fontenc}
\usepackage[utf8]{inputenc}
\usepackage{mathtools}
\usepackage{amsmath}
\usepackage{amssymb}

\makeatletter
%%%%%%%%%%%%%%%%%%%%%%%%%%%%%% Textclass specific LaTeX commands.
\numberwithin{equation}{section}
\numberwithin{figure}{section}

%%%%%%%%%%%%%%%%%%%%%%%%%%%%%% User specified LaTeX commands.
\usepackage[left=1.4cm, right=1.5cm, top=1.5cm, bottom=1.5cm]{geometry}

\makeatother

\usepackage{babel}
\begin{document}

\section{Unconditional Non-linear denoiser}

Random-feature model imposes a non-linear map on the input of the
denoiser $\boldsymbol{y}$, however the denoiser output is still linear
in this new non-linear feature map; let $\psi$ be a non-linear function,
i.e. sigmoid, softmax, etc, $\Theta$ be a $\mathbb{R}^{d\times k}$
random matrix, e.g Gaussian entries, $\epsilon$ be a $\mathbb{R}^{d}$
random bias; all are fixed i.e. we're not learning these here, as
opposed to the way a full neural network does. Then:
\[
\phi(\boldsymbol{y})=\omega(\Theta\boldsymbol{y}+\epsilon)
\]

And the denoiser w.r.t this new non-linear feature map is still:
\[
\boldsymbol{D}_{\theta}(\phi(\boldsymbol{y});\sigma,\boldsymbol{U})=\boldsymbol{W}_{\sigma}\phi(\boldsymbol{y})+\boldsymbol{b}_{\sigma}
\]

where $\boldsymbol{W}_{\sigma}$ and $\boldsymbol{b}_{\sigma}$ are
learned. We have:
\[
\mathcal{L}_{\sigma}=\mathbb{E}_{\boldsymbol{x_{0}}\sim p_{0},\boldsymbol{Z}\sim\mathcal{N}(0,\boldsymbol{I})}\|\boldsymbol{W}_{\sigma}\phi(\boldsymbol{x_{0}}+\sigma\boldsymbol{Z})+\boldsymbol{b}_{\sigma}-\boldsymbol{x_{0}}\|^{2}_{2}
\]

Note that the denoiser is linear w.r.t the feature map so most of
our derivations from the linear case still applies. Setting derivative
w.r.t $\boldsymbol{b}_{\sigma}$ to $\boldsymbol{0}$:
\begin{align*}
\nabla_{\boldsymbol{b}_{\sigma}}\mathcal{L}_{\sigma} & =\mathbb{E}_{\boldsymbol{x_{0}}\sim p_{0},\boldsymbol{Z}\sim\mathcal{N}(0,\boldsymbol{I})}\big[2(\boldsymbol{W}_{\sigma}\phi(\boldsymbol{x_{0}}+\sigma\boldsymbol{Z})+\boldsymbol{b}_{\sigma}-\boldsymbol{x_{0}})\big]\\
0 & =2(\boldsymbol{W}_{\sigma}\mu_{\phi}+\boldsymbol{b}_{\sigma}-\mu_{\boldsymbol{p_{0}}})
\end{align*}

where $\mu_{\phi}=\mathbb{E}_{\boldsymbol{x_{0}}\sim p_{0},\boldsymbol{Z}\sim\mathcal{N}(0,\boldsymbol{I})}[\phi(\boldsymbol{x_{0}}+\sigma\boldsymbol{Z})]$.
So:
\[
\boldsymbol{b}^{*}_{\sigma}=\mu_{\boldsymbol{p_{0}}}-\boldsymbol{W}_{\sigma}\mu_{\phi}
\]

Similarly:
\begin{align*}
 & \nabla_{\boldsymbol{W}_{\sigma}}\mathbb{E}_{\boldsymbol{x_{0}}\sim p_{0},\boldsymbol{Z}\sim\mathcal{N}(0,\boldsymbol{I})}\|\boldsymbol{W}_{\sigma}\phi(\boldsymbol{x_{0}}+\sigma\boldsymbol{Z})+\boldsymbol{b}_{\sigma}-\boldsymbol{x_{0}}\|^{2}_{2}\\
 & =\mathbb{E}_{\boldsymbol{x_{0}}\sim p_{0},\boldsymbol{Z}\sim\mathcal{N}(0,\boldsymbol{I})}\big[2(\boldsymbol{W}_{\sigma}\phi(\boldsymbol{x_{0}}+\sigma\boldsymbol{Z})+\boldsymbol{b}_{\sigma}-\boldsymbol{x_{0}})\phi(\boldsymbol{x_{0}}+\sigma\boldsymbol{Z})^{\top}\big]\\
 & =2(\boldsymbol{W}_{\sigma}(\boldsymbol{\Sigma}_{\phi}+\mu_{\phi}\mu^{\top}_{\phi})+(\mu_{\boldsymbol{p_{0}}}-\boldsymbol{W}_{\sigma}\mu_{\phi})\mu_{\phi}-\mathbb{E}[\boldsymbol{x_{0}}\phi(\boldsymbol{x_{0}}+\sigma\boldsymbol{Z})^{\top}])\\
0 & =2(\boldsymbol{W}_{\sigma}(\boldsymbol{\Sigma}_{\phi})+\text{Cov}(\boldsymbol{x_{0}},\phi))
\end{align*}

where $\boldsymbol{\Sigma}_{\phi}=\text{Cov}(\phi(\boldsymbol{x_{0}}+\sigma\boldsymbol{Z}))$.
So $\boldsymbol{W}^{*}_{\sigma}=\text{Cov}(\boldsymbol{x_{0}},\phi)\boldsymbol{\Sigma}^{-1}_{\phi}$.
Then:
\begin{align*}
\mathcal{L}_{\sigma} & =\mathbb{E}_{\boldsymbol{x_{0}}\sim p_{0},\boldsymbol{Z}\sim\mathcal{N}(0,\boldsymbol{I})}\|\boldsymbol{W}_{\sigma}(\phi(\boldsymbol{x_{0}}+\sigma\boldsymbol{Z})-\mu_{\phi})-(\boldsymbol{x_{0}}-\mu_{\boldsymbol{p_{0}}})\|^{2}_{2}\\
 & =\text{Tr}(\mathbb{E}_{\boldsymbol{x_{0}}\sim p_{0},\boldsymbol{Z}\sim\mathcal{N}(0,\boldsymbol{I})}[\boldsymbol{W}_{\sigma}(\phi(\boldsymbol{x_{0}}+\sigma\boldsymbol{Z})-\mu_{\phi})(\phi(\boldsymbol{x_{0}}+\sigma\boldsymbol{Z})-\mu_{\phi})^{\top}\boldsymbol{W}^{\top}_{\sigma}]\\
 & -2\mathbb{E}_{\boldsymbol{x_{0}}\sim p_{0},\boldsymbol{Z}\sim\mathcal{N}(0,\boldsymbol{I})}[\boldsymbol{W}_{\sigma}(\phi(\boldsymbol{x_{0}}+\sigma\boldsymbol{Z})-\mu_{\phi})(\boldsymbol{x_{0}}-\mu_{\boldsymbol{p_{0}}})^{\top}]+\mathbb{E}_{\boldsymbol{x_{0}}\sim p_{0}}[(\boldsymbol{x_{0}}-\mu_{\boldsymbol{p_{0}}})(\boldsymbol{x_{0}}-\mu_{\boldsymbol{p_{0}}})^{\top}])\\
 & =\text{Tr}(\boldsymbol{W}_{\sigma}\boldsymbol{\Sigma}_{\phi(\boldsymbol{x_{0}}+\sigma\boldsymbol{Z})-\mu_{\phi}}\boldsymbol{W}^{\top}_{\sigma}-2\boldsymbol{W}_{\sigma}\text{Cov}(\boldsymbol{x_{0}},\phi)+\boldsymbol{\Sigma}_{\boldsymbol{x_{0}}-\mu_{\boldsymbol{p_{0}}}})\\
 & =\text{Tr}(\boldsymbol{W}_{\sigma}\boldsymbol{\Sigma}_{\phi}\boldsymbol{W}^{\top}_{\sigma}-2\boldsymbol{W}_{\sigma}\text{Cov}(\boldsymbol{x_{0}},\phi)+\boldsymbol{\Sigma}_{\boldsymbol{p_{0}}})\\
 & =\text{Tr}(\text{Cov}(\boldsymbol{x_{0}},\phi)\boldsymbol{\Sigma}^{-1}_{\phi}\boldsymbol{\Sigma}_{\phi}(\text{Cov}(\boldsymbol{x_{0}},\phi)\boldsymbol{\Sigma}^{-1}_{\phi})^{\top}-2\text{Cov}(\boldsymbol{x_{0}},\phi)\boldsymbol{\Sigma}^{-1}_{\phi}\text{Cov}(\boldsymbol{x_{0}},\phi)+\boldsymbol{\Sigma}_{\boldsymbol{p_{0}}})\\
 & =\text{Tr}(\boldsymbol{\Sigma}_{\boldsymbol{p_{0}}}-\text{Cov}(\boldsymbol{x_{0}},\phi)\boldsymbol{\Sigma}^{-1}_{\phi}\text{Cov}(\phi,\boldsymbol{x_{0}}))
\end{align*}

where $\text{Cov}(\boldsymbol{x_{0}},\phi)=\mathbb{E}_{\boldsymbol{x_{0}}\sim p_{0},\boldsymbol{Z}\sim\mathcal{N}(0,\boldsymbol{I})}[(\phi(\boldsymbol{x_{0}}+\sigma\boldsymbol{Z})-\mu_{\phi})(\boldsymbol{x_{0}}-\mu_{\boldsymbol{p_{0}}})^{\top}]$

Now if $\phi$ is Gaussian square-integrable w.r.t the Gaussian measure,
we can use Hermite polynomials to approximate it. Recall:
\[
\text{He}_{n}=(-1)^{n}e^{x^{2}/2}\frac{d^{n}}{dx^{n}}(e^{-x^{2}/2})
\]

The first few look like:
\begin{align*}
\text{He}_{0} & =1\\
\text{He}_{1} & =x-1\\
\text{He}_{2} & =x^{2}-1\\
\text{He}_{3} & =x^{3}-3x\\
\text{He}_{4} & =x^{4}-6x^{2}+3\\
 & \vdots
\end{align*}

We can show that:
\[
\text{He}_{n+1}(x)=x\text{He}_{n}(x)-n\text{He}_{n-1}(x)
\]

These polynomials are the orthogonal basis of $\mathcal{P}(x)$, the
space of all polynomials in $x$; their orthogonality is under the
inner-product of the space which is with respect to the Gaussian measure
$\langle\cdot,\cdot\rangle_{\Phi}$, given as $\langle p,q\rangle_{\varphi}=\int p(x)q(x)\varphi(x)dx$
where $\varphi(x)=\frac{e^{-x^{2}/2}}{\sqrt{2\pi}}.$ In other words,
$\langle p,q\rangle_{\varphi}=\mathbb{E}_{Z\sim\mathcal{N}(0,1)}[p(Z)q(Z)]$.
Then:
\begin{align*}
f(x) & =\sum_{n}\frac{1}{n!}\langle f(x),\text{He}_{n}(x)\rangle_{\varphi}\text{He}_{n}(x)
\end{align*}

Denote $\frac{1}{n!}\langle f(x),\text{He}_{n}(x)\rangle_{\varphi}=c_{n}(x)$

First consider $\mu_{\phi}$:
\begin{align*}
\mathbb{E}_{\boldsymbol{x_{0}}\sim p_{0},\boldsymbol{Z}\sim\mathcal{N}(0,\boldsymbol{I})}[\phi(\boldsymbol{x_{0}}+\sigma\boldsymbol{Z})] & =\mathbb{E}_{\boldsymbol{x_{0}}\sim p_{0}}[\mathbb{E}_{\boldsymbol{Z}\sim\mathcal{N}(0,\boldsymbol{I})}[\phi(\boldsymbol{x_{0}}+\sigma\boldsymbol{Z})]]\\
 & =\mathbb{E}_{\boldsymbol{x_{0}}\sim p_{0}}[\mathbb{E}_{\boldsymbol{Z}\sim\mathcal{N}(0,\boldsymbol{I})}[\sum_{n}c_{n}(\boldsymbol{x_{0}}+\sigma\boldsymbol{Z})\text{He}_{n}(\boldsymbol{x_{0}}+\sigma\boldsymbol{Z})]]\\
 & =\mathbb{E}_{\boldsymbol{x_{0}}\sim p_{0}}[\sum_{n}c_{n}\mathbb{E}_{\boldsymbol{Z}\sim\mathcal{N}(0,\boldsymbol{I})}[\text{He}_{n}(\boldsymbol{x_{0}}+\sigma\boldsymbol{Z})]]\\
 & =\mathbb{E}_{\boldsymbol{x_{0}}\sim p_{0}}[\sum_{n}\langle\text{He}_{0},c_{n}(\boldsymbol{x_{0}}+\sigma\boldsymbol{Z})\text{He}_{n}\rangle_{\varphi}]
\end{align*}

We want to expand the closed-form of $\mathcal{L}_{\sigma}$. Here
if we Hermite-expand on the non-linear $\phi(\boldsymbol{x_{0}}+\sigma\boldsymbol{Z})$
then our Hermite polynomials would look like $\text{He}_{n}(\boldsymbol{x_{0}}+\sigma\boldsymbol{Z})$
and coefficients like $c_{n}(\boldsymbol{x_{0}}+\sigma\boldsymbol{Z})$,
which is not ideal for simplifying any further since we will be dealing
with $\boldsymbol{x_{0}}+\sigma\boldsymbol{Z}$ which is not Gaussian
unless $\boldsymbol{x_{0}}$ is Gaussian. So first condition on $\boldsymbol{x_{0}}$
first and then expand. Consider the case $\phi$ is a component-wise
non-linear function; then we can get around $\phi(\boldsymbol{y})=\omega(\Theta\boldsymbol{y}+\xi)$
by building it row by row. Let $\theta_{j}$ be the $j^{\text{th}}$
row of $\Theta$, and conditional on $\boldsymbol{x_{0}}$:
\begin{align*}
\phi_{j}(\boldsymbol{x_{0}}+\sigma\boldsymbol{Z}) & =\omega(\theta_{j}(\boldsymbol{x_{0}}+\sigma\boldsymbol{Z})+\epsilon_{j})=\omega(\theta^{\top}_{j}\boldsymbol{x_{0}}+\sigma\theta^{\top}_{j}\boldsymbol{Z}+\epsilon_{j})\coloneqq\omega(M_{j}(\boldsymbol{x_{0}})+s_{j}\xi);\\
\text{where} & \qquad\qquad\:s_{j}=\sigma\|\theta_{j}\|,\quad\xi=\frac{\theta^{\top}_{j}\boldsymbol{Z}}{\|\theta_{j}\|}\sim\mathcal{N}(0,1),\quad M_{j}(\boldsymbol{x_{0}})=\theta^{\top}_{j}\boldsymbol{x_{0}}+\epsilon_{j}
\end{align*}

Conditional on $\boldsymbol{x_{0}},$ $M_{j}(\boldsymbol{x_{0}})$
is just a constant; and $s_{j}$ is also a constant. Then $\omega$
really is a function of $\xi$; let $f(\xi):\xi\mapsto\omega(M_{j}(\boldsymbol{x_{0}})+s_{j}\xi)$.
Now we Hermite-expand $f(\xi)$:
\begin{align*}
\omega(M_{j}(\boldsymbol{x_{0}})+s_{j}\xi) & =\sum_{n=0}c_{n}(M_{j}(\boldsymbol{x_{0}}),s_{j})\text{He}_{n}(\xi),\quad\quad\quad c_{n}(M_{j}(\boldsymbol{x_{0}}),s_{j})=\frac{1}{n!}\langle f(\xi),\text{He}_{n}(\xi)\rangle_{\varphi}=\mathbb{E}_{\xi\sim\mathcal{N}(0,1)}[f(\xi)\text{He}_{n}(\xi)]
\end{align*}

Note that $c_{n}(M_{j}(\boldsymbol{x_{0}}),s_{j})$ is a constant
w.r.t $\xi$ since $\xi$ would be integrated out. Also that $\mathbb{E}_{\xi\sim\mathcal{N}(0,1)}[\text{He}_{i}(\xi)\text{He}_{j}(\xi)]=0$
by orthogonality. Then by Tower property:
\begin{align*}
\mathbb{E}[\phi_{j}(\boldsymbol{x_{0}}+\sigma\boldsymbol{Z})] & =\mathbb{E}_{\boldsymbol{x_{0}}\sim\boldsymbol{p_{0}}}[\mathbb{E}[\phi_{j}(\boldsymbol{x_{0}}+\sigma\boldsymbol{Z})\mid\boldsymbol{x_{0}}]]\\
 & =\mathbb{E}_{\boldsymbol{x_{0}}\sim\boldsymbol{p_{0}}}[\mathbb{E}_{\xi\sim\mathcal{N}(0,1)}[\sum_{n=0}c_{n}(M_{j}(\boldsymbol{x_{0}}),s_{j})\text{He}_{n}(\xi)]]\\
 & =\mathbb{E}_{\boldsymbol{x_{0}}\sim\boldsymbol{p_{0}}}[\sum_{n=0}c_{n}(M_{j}(\boldsymbol{x_{0}}),s_{j})\mathbb{E}_{\xi\sim\mathcal{N}(0,1)}[\text{He}_{n}(\xi)\text{He}_{0}(\xi)]]\quad\quad(\text{since }\text{He}_{0}=1)\\
 & =\mathbb{E}_{\boldsymbol{x_{0}}\sim\boldsymbol{p_{0}}}[c_{0}(M_{j}(\boldsymbol{x_{0}}),s_{j})]
\end{align*}

Rewrite $c_{0}(M_{j}(\boldsymbol{x_{0}}),s_{j})\coloneqq g_{j}(\boldsymbol{x_{0}})$
since $g_{j}$ varies with $\boldsymbol{x_{0}}$:
\begin{align*}
g_{j}(\boldsymbol{x_{0}}) & =\mathbb{E}_{\xi\sim\mathcal{N}(0,1)}[f(\xi)\text{He}_{0}(\xi)]\\
 & =\mathbb{E}_{\xi\sim\mathcal{N}(0,1)}[\omega(M_{j}(\boldsymbol{x_{0}})+s_{j}\xi)]\\
 & =\int\omega(M_{j}(\boldsymbol{x_{0}})+s_{j}\xi)\frac{e^{-\xi^{2}/2}}{\sqrt{2\pi}}d\xi
\end{align*}

Let $u=s_{j}\xi$, so $\xi=\frac{u}{s_{j}}$ and $d\xi=\frac{du}{s_{j}}$:
\begin{align*}
g_{j}(\boldsymbol{x_{0}}) & =\int\omega(M_{j}(\boldsymbol{x_{0}})+u)\frac{e^{-u^{2}/2s^{2}_{j}}}{s_{j}\sqrt{2\pi}}du\\
\end{align*}

This integral is the convolution of function $\omega$ and the Gaussian
$\mathcal{N}(0,s^{2}_{j})$, evaluated at the local point $M_{j}(\boldsymbol{x_{0}})$,
i.e.
\begin{align*}
g_{j}(\boldsymbol{x_{0}}) & =(\omega*\mathcal{N}_{s^{2}_{j}})(M_{j}(\boldsymbol{x_{0}}))
\end{align*}

The interpretation is that this is a function $\omega$ blurred by
a Gaussian kernel of width $s_{j}$, then evaluated at the point $M_{j}(\boldsymbol{x_{0}})$.
Now consider $\text{Cov}(\boldsymbol{x_{0}},\phi)_{ij}$ at row $i^{\text{th}}$
and column $j^{\text{th}}$:
\begin{align*}
\text{Cov}(\boldsymbol{x_{0}},\phi)_{ij} & =\mathbb{E}_{\boldsymbol{x_{0}}\sim p_{0},\boldsymbol{Z}\sim\mathcal{N}(0,\boldsymbol{I})}[(\boldsymbol{x}_{\boldsymbol{0},i}-\mu_{\boldsymbol{p_{0}},i})(\phi(\boldsymbol{x_{0}}+\sigma\boldsymbol{Z}))_{j}]\\
 & =\mathbb{E}_{\boldsymbol{x_{0}}\sim p_{0}}[\mathbb{E}_{\boldsymbol{Z}\sim\mathcal{N}(0,\boldsymbol{I})}[(\boldsymbol{x}_{\boldsymbol{0},i}-\mu_{\boldsymbol{p_{0}},i})(\phi(\boldsymbol{x_{0}}+\sigma\boldsymbol{Z}))_{j}\mid\boldsymbol{x_{0}}]]\\
 & =\mathbb{E}_{\boldsymbol{x_{0}}\sim p_{0}}[(\boldsymbol{x}_{\boldsymbol{0},i}-\mu_{\boldsymbol{p_{0}},i})\mathbb{E}_{\boldsymbol{Z}\sim\mathcal{N}(0,\boldsymbol{I})}[(\phi(\boldsymbol{x_{0}}+\sigma\boldsymbol{Z}))_{j}\mid\boldsymbol{x_{0}}]]\\
 & =\mathbb{E}_{\boldsymbol{x_{0}}\sim p_{0}}[(\boldsymbol{x}_{\boldsymbol{0},i}-\mu_{\boldsymbol{p_{0}},i})g_{j}(\boldsymbol{x_{0}})]
\end{align*}

Note that $g_{j}(\boldsymbol{x_{0}})=c_{0}(M_{j}(\boldsymbol{x_{0}}),s_{j})$
where $M_{j}(\boldsymbol{x_{0}})=\Theta^{\top}_{j}\boldsymbol{x_{0}}+\epsilon_{j}$
so really $g_{j}(\boldsymbol{x_{0}})$ only depends on $u_{j}\coloneqq\Theta^{\top}_{j}\boldsymbol{x_{0}}+\epsilon_{j}$,
so rewrite $g_{j}(\boldsymbol{x_{0}})=\psi(\Theta^{\top}_{j}\boldsymbol{x_{0}}+\epsilon_{j})=\psi(u_{j})$
where $\psi(u)=c_{0}(u,s_{j})$. Now by Tower property:
\begin{align*}
\text{Cov}(\boldsymbol{x_{0}},\phi)_{ij} & =\mathbb{E}_{\boldsymbol{x_{0}}\sim p_{0}}[(\boldsymbol{x}_{\boldsymbol{0},i}-\mu_{\boldsymbol{p_{0}},i})g_{j}(\boldsymbol{x_{0}})]\\
 & =\mathbb{E}_{u_{j}}[\mathbb{E}_{\boldsymbol{x_{0}}\sim p_{0}}[(\boldsymbol{x}_{\boldsymbol{0},i}-\mu_{\boldsymbol{p_{0}},i})g_{j}(\boldsymbol{x_{0}})\mid u_{j}]]\\
 & =\mathbb{E}_{u_{j}}[\mathbb{E}_{\boldsymbol{x_{0}}\sim p_{0}}[(\boldsymbol{x}_{\boldsymbol{0},i}-\mu_{\boldsymbol{p_{0}},i})g_{j}(\boldsymbol{x_{0}})\mid u_{j}]]\\
 & =\mathbb{E}_{u_{j}}[g_{j}(u_{j})\mathbb{E}_{\boldsymbol{x_{0}}\sim p_{0}}[(\boldsymbol{x}_{\boldsymbol{0},i}-\mu_{\boldsymbol{p_{0}},i})\mid u_{j}]]
\end{align*}

Define $r_{i}(u_{j})=\mathbb{E}_{\boldsymbol{x_{0}}\sim p_{0}}[(\boldsymbol{x}_{\boldsymbol{0},i}-\mu_{\boldsymbol{p_{0}},i})\mid u_{j}]]$,
which is really a regression coefficient of $\boldsymbol{x_{0}}$
onto its own linear projection. All this do is turning a multi-dimensional
integral $\mathbb{E}_{\boldsymbol{x_{0}}\sim p_{0}}$ into an 1-D
integral $\mathbb{E}_{u_{j}}$ .

Now consider $\boldsymbol{\Sigma}_{\phi}$; recall that $\omega$
is component-wise so we're allowed to break $\phi$ down into its
entries. 
\begin{align*}
\boldsymbol{\Sigma}_{\phi,ij} & =\text{Cov}(\phi_{i},\phi_{j})=\mathbb{E}_{\boldsymbol{x_{0}},\boldsymbol{Z}}[\phi_{i}\phi_{j}]-\mathbb{E}_{\boldsymbol{x_{0}},\boldsymbol{Z}}[\phi_{i}]\mathbb{E}_{\boldsymbol{x_{0}},\boldsymbol{Z}}[\phi_{j}]
\end{align*}

Consider the first expectation term; by the Tower property:
\begin{align*}
\mathbb{E}_{\boldsymbol{x_{0}},\boldsymbol{Z}}[\phi_{i}\phi_{j}] & =\mathbb{E}_{\boldsymbol{x_{0}}}[\mathbb{E}_{\boldsymbol{Z}}[\phi_{i}(\boldsymbol{x_{0}}+\sigma\boldsymbol{Z})\phi_{j}(\boldsymbol{x_{0}}+\sigma\boldsymbol{Z})\mid\boldsymbol{x_{0}}]]\\
 & =\mathbb{E}_{\boldsymbol{x_{0}}}\big[\text{Cov}(\phi_{i},\phi_{j}\mid\boldsymbol{x_{0}})+\mathbb{E}_{\boldsymbol{Z}}[\phi_{j}(\boldsymbol{x_{0}}+\sigma\boldsymbol{Z})\mid\boldsymbol{x_{0}}]\mathbb{E}_{\boldsymbol{Z}}[\phi_{j}(\boldsymbol{x_{0}}+\sigma\boldsymbol{Z})\mid\boldsymbol{x_{0}}]\big]\\
 & =\mathbb{E}_{\boldsymbol{x_{0}}}[\text{Cov}(\phi_{i},\phi_{j}\mid\boldsymbol{x_{0}})]+\mathbb{E}_{\boldsymbol{x_{0}}}[g_{i}(\boldsymbol{x_{0}})g_{j}(\boldsymbol{x_{0}})]
\end{align*}

Now:
\begin{align*}
\mathbb{E}_{\boldsymbol{x_{0}},\boldsymbol{Z}}[\phi_{i}]\mathbb{E}_{\boldsymbol{x_{0}},\boldsymbol{Z}}[\phi_{j}] & =\mathbb{E}_{\boldsymbol{x_{0}}}[\mathbb{E}_{\boldsymbol{Z}}[\phi_{i}\mid\boldsymbol{x_{0}}]]\mathbb{E}_{\boldsymbol{x_{0}}}[\mathbb{E}_{\boldsymbol{Z}}[\phi_{i}\mid\boldsymbol{x_{0}}]]\\
 & =\mathbb{E}_{\boldsymbol{x_{0}}}[g_{i}(\boldsymbol{x_{0}})]\mathbb{E}_{\boldsymbol{x_{0}}}[g_{j}(\boldsymbol{x_{0}})]
\end{align*}

So:
\begin{align*}
\boldsymbol{\Sigma}_{\phi,ij} & =\mathbb{E}_{\boldsymbol{x_{0}},\boldsymbol{Z}}[\phi_{i}\phi_{j}]-\mathbb{E}_{\boldsymbol{x_{0}},\boldsymbol{Z}}[\phi_{i}]\mathbb{E}_{\boldsymbol{x_{0}},\boldsymbol{Z}}[\phi_{j}]\\
 & =\mathbb{E}_{\boldsymbol{x_{0}}}[\text{Cov}_{\boldsymbol{Z}}(\phi_{i},\phi_{j}\mid\boldsymbol{x_{0}})]+\mathbb{E}_{\boldsymbol{x_{0}}}[g_{i}(\boldsymbol{x_{0}})g_{j}(\boldsymbol{x_{0}})]-\mathbb{E}_{\boldsymbol{x_{0}}}[g_{i}(\boldsymbol{x_{0}})]\mathbb{E}_{\boldsymbol{x_{0}}}[g_{j}(\boldsymbol{x_{0}})]\\
 & =\mathbb{E}_{\boldsymbol{x_{0}}}[\text{Cov}_{\boldsymbol{Z}}(\phi_{i},\phi_{j}\mid\boldsymbol{x_{0}})]+\text{Cov}_{\boldsymbol{x_{0}}}(g_{i}(\boldsymbol{x_{0}}),g_{j}(\boldsymbol{x_{0}}))
\end{align*}

Consider the term $\mathbb{E}_{\boldsymbol{x_{0}}}[\text{Cov}(\phi_{i},\phi_{j}\mid\boldsymbol{x_{0}})]$:
\begin{align*}
\mathbb{E}_{\boldsymbol{x_{0}}}[\text{Cov}_{\boldsymbol{Z}}(\phi_{i},\phi_{j}\mid\boldsymbol{x_{0}})] & =\mathbb{E}_{\boldsymbol{x_{0}}}[\mathbb{E}_{\boldsymbol{Z}}[\phi_{i}\phi_{j}\mid\boldsymbol{x_{0}}]-\mathbb{E}_{\boldsymbol{Z}}[\phi_{i}\mid\boldsymbol{x_{0}}]\mathbb{E}_{\boldsymbol{Z}}[\phi_{j}\mid\boldsymbol{x_{0}}]]\\
 & =\mathbb{E}_{\boldsymbol{x_{0}}}[\mathbb{E}_{\boldsymbol{Z}}[\phi_{i}\phi_{j}\mid\boldsymbol{x_{0}}]-c_{0}(M_{i}(\boldsymbol{x_{0}}),s_{i})c_{0}(M_{j}(\boldsymbol{x_{0}}),s_{j})]
\end{align*}

Recall Hermite-expansion of $\phi$ conditional on $\boldsymbol{x_{0}}$:
\begin{align*}
\phi_{i} & =\sum_{n=0}c_{n}(M_{i}(\boldsymbol{x_{0}}),s_{i})\text{He}_{n}(\xi_{i}),\quad\phi_{j}=\sum_{m=0}c_{m}(M_{j}(\boldsymbol{x_{0}}),s_{j})\text{He}_{m}(\xi_{j});\hspace*{\fill}
\end{align*}

So:
\begin{align*}
\mathbb{E}_{\boldsymbol{Z}}[\phi_{i}\phi_{j}\mid\boldsymbol{x_{0}}] & =\mathbb{E}_{\boldsymbol{Z}}\big[\sum_{n=0}c_{n}(M_{i}(\boldsymbol{x_{0}}),s_{i})\text{He}_{n}(\xi_{i})\sum_{m=0}c_{m}(M_{j}(\boldsymbol{x_{0}}),s_{j})\text{He}_{m}(\xi_{j})\big]\\
 & =\sum_{n,m}c_{n}(M_{i}(\boldsymbol{x_{0}}),s_{i})c_{m}(M_{j}(\boldsymbol{x_{0}}),s_{j})\mathbb{E}\big[\text{He}_{n}(\xi_{i})\text{He}_{m}(\xi_{j})\big]
\end{align*}

Using the Mehler product identity for correlated Gaussian:
\begin{align*}
\mathbb{E}\big[\text{He}_{n}(\xi_{i})\text{He}_{m}(\xi_{j})\big] & =n!\rho^{n}_{ij}\delta_{nm}
\end{align*}

where $\rho_{ij}=\text{Corr}(\xi_{i},\xi_{j})=\frac{\theta^{\top}_{i}\theta_{j}}{\|\theta_{i}\|\|\theta_{j}\|}$
and $\delta_{nm}$ is the Kronecker delta. So:
\begin{align*}
\mathbb{E}_{\boldsymbol{Z}}[\phi_{i}\phi_{j}\mid\boldsymbol{x_{0}}] & =\sum_{n}\sum_{m}n!c_{n}(M_{i}(\boldsymbol{x_{0}}),s_{i})c_{m}(M_{j}(\boldsymbol{x_{0}}),s_{j})\rho^{n}_{ij}\delta_{nm}\\
 & =\sum_{n}n!\rho^{n}_{ij}c_{n}(M_{i}(\boldsymbol{x_{0}}),s_{i})\sum_{m}c_{m}(M_{j}(\boldsymbol{x_{0}}),s_{j})\delta_{nm}\\
 & =\sum_{n}n!\rho^{n}_{ij}c_{n}(M_{i}(\boldsymbol{x_{0}}),s_{i})c_{n}(M_{j}(\boldsymbol{x_{0}}),s_{j})
\end{align*}

Therefore:
\begin{align*}
\mathbb{E}_{\boldsymbol{x_{0}}}[\text{Cov}_{\boldsymbol{Z}}(\phi_{i},\phi_{j}\mid\boldsymbol{x_{0}})] & =\mathbb{E}_{\boldsymbol{x_{0}}}[\sum_{n}n!\rho^{n}_{ij}c_{n}(M_{i}(\boldsymbol{x_{0}}),s_{i})c_{n}(M_{j}(\boldsymbol{x_{0}}),s_{j})-c_{0}(M_{i}(\boldsymbol{x_{0}}),s_{i})c_{0}(M_{j}(\boldsymbol{x_{0}}),s_{j})]\\
 & =\mathbb{E}_{\boldsymbol{x_{0}}}[\sum_{n\geq1}n!\rho^{n}_{ij}c_{n}(M_{i}(\boldsymbol{x_{0}}),s_{i})c_{n}(M_{j}(\boldsymbol{x_{0}}),s_{j})]
\end{align*}

So:
\[
\boldsymbol{\Sigma}_{\phi,ij}=\text{Cov}_{\boldsymbol{x_{0}}}(g_{i}(\boldsymbol{x_{0}}),g_{j}(\boldsymbol{x_{0}}))+\mathbb{E}_{\boldsymbol{x_{0}}}[\sum_{n\geq1}n!\rho^{n}_{ij}c_{n}(M_{i}(\boldsymbol{x_{0}}),s_{i})c_{n}(M_{j}(\boldsymbol{x_{0}}),s_{j})]
\]

To summarize:
\[
\mathcal{L}_{\sigma}=\text{Tr}(\boldsymbol{\Sigma}_{\boldsymbol{p_{0}}}-\text{Cov}(\boldsymbol{x_{0}},\phi)\boldsymbol{\Sigma}^{-1}_{\phi}\text{Cov}(\phi,\boldsymbol{x_{0}}))
\]

where the entries of $\boldsymbol{\Sigma}_{\phi}$ is defined as above
and the entries of $\text{Cov}(\boldsymbol{x_{0}},\phi)$ is 
\[
\text{Cov}(\boldsymbol{x_{0}},\phi)_{ij}=\mathbb{E}_{u_{j}}[g_{j}(u_{j})\mathbb{E}_{\boldsymbol{x_{0}}\sim p_{0}}[(\boldsymbol{x}_{\boldsymbol{0},i}-\mu_{\boldsymbol{p_{0}},i})\mid u_{j}]]
\]


\section{Conditional Non-linear denoiser}

In the conditional case, if we treat $(\boldsymbol{y},\boldsymbol{U})$
as a single joint input vector where $\phi^{\boldsymbol{U}}(\boldsymbol{y},\boldsymbol{U})=\omega(\Theta\boldsymbol{y}+\Gamma\boldsymbol{U}+\epsilon)$
is a random feature map then the loss has the same form as before
except that:
\begin{align*}
\mu^{\boldsymbol{U}}_{\phi} & =\mathbb{E}[\phi^{\boldsymbol{U}}(\boldsymbol{x_{0}}+\sigma\boldsymbol{Z},\boldsymbol{U})]\\
\boldsymbol{\Sigma}^{\boldsymbol{U}}_{\phi} & =\text{Cov}(\phi(\boldsymbol{x_{0}}+\sigma\boldsymbol{Z},\boldsymbol{U}))\\
\text{Cov}(\boldsymbol{x_{0}},\phi^{\boldsymbol{U}}) & =\mathbb{E}[(\boldsymbol{x_{0}}-\mu_{\boldsymbol{p_{0}}})(\phi^{\boldsymbol{U}}(\boldsymbol{x_{0}}+\sigma\boldsymbol{Z},\boldsymbol{U})-\mu^{\boldsymbol{U}}_{\phi})^{\top}]
\end{align*}

and so:
\[
\mathcal{L}_{\sigma,\boldsymbol{U}}=\text{Tr}(\boldsymbol{\Sigma}_{\boldsymbol{p_{0}}}-\text{Cov}(\boldsymbol{x_{0}},\phi^{\boldsymbol{U}})(\boldsymbol{\Sigma}^{\boldsymbol{U}}_{\phi})^{-1}\text{Cov}(\phi^{\boldsymbol{U}},\boldsymbol{x_{0}}))
\]

We want to expand this more. Carrying over from the unconditional
case, let $\theta_{j}$ be the $j^{\text{th}}$ row of $\Theta$,
$\gamma_{j}$ be the $j^{\text{th}}$ row of $\Gamma$, and conditional
on $\boldsymbol{x_{0}}$:
\begin{align*}
\phi_{j}(\boldsymbol{x_{0}}+\sigma\boldsymbol{Z}) & =\omega(\theta_{j}(\boldsymbol{x_{0}}+\sigma\boldsymbol{Z})+\gamma_{j}\boldsymbol{U}+\epsilon_{j})=\omega(\theta^{\top}_{j}\boldsymbol{x_{0}}+\sigma\theta^{\top}_{j}\boldsymbol{Z}+\gamma^{\top}_{j}\boldsymbol{U}+\epsilon_{j})\coloneqq\omega(M^{\boldsymbol{U}}_{j}(\boldsymbol{x_{0}})+s_{j}\xi);\\
\text{where} & \qquad\qquad\:s_{j}=\sigma\|\theta_{j}\|,\quad\xi=\frac{\theta^{\top}_{j}\boldsymbol{Z}}{\|\theta_{j}\|}\sim\mathcal{N}(0,1),\quad M^{\boldsymbol{U}}_{j}(\boldsymbol{x_{0}})=\theta^{\top}_{j}\boldsymbol{x_{0}}+\gamma_{j}\boldsymbol{U}+\epsilon_{j}
\end{align*}

We're given $\boldsymbol{U}$ and conditional on $\boldsymbol{x_{0}}$,
$\omega$ really is a function of $\xi$; let $f(\xi):\xi\mapsto\omega(M^{\boldsymbol{U}}_{j}(\boldsymbol{x_{0}})+s_{j}\xi)$.
Hermite-expand $f(\xi)$:
\begin{align*}
\omega(M^{\boldsymbol{U}}_{j}(\boldsymbol{x_{0}})+s_{j}\xi) & =\sum_{n=0}c_{n}(M^{\boldsymbol{U}}_{j}(\boldsymbol{x_{0}}),s_{j})\text{He}_{n}(\xi),\quad\quad\quad c_{n}(M^{\boldsymbol{U}}_{j}(\boldsymbol{x_{0}}),s_{j})=\frac{1}{n!}\langle f(\xi),\text{He}_{n}(\xi)\rangle_{\varphi}=\mathbb{E}_{\xi\sim\mathcal{N}(0,1)}[f(\xi)\text{He}_{n}(\xi)]
\end{align*}

Note that $c_{n}(M^{\boldsymbol{U}}_{j}(\boldsymbol{x_{0}}),s_{j})$
is a constant w.r.t $\xi$ since $\xi$ would be integrated out. Also
that $\mathbb{E}_{\xi\sim\mathcal{N}(0,1)}[\text{He}_{i}(\xi)\text{He}_{j}(\xi)]=0$
by orthogonality. Then by Tower property:
\begin{align*}
\mathbb{E}[\phi^{\boldsymbol{U}}_{j}(\boldsymbol{x_{0}}+\sigma\boldsymbol{Z},\boldsymbol{U})] & =\mathbb{E}_{\boldsymbol{x_{0}}\sim\boldsymbol{p_{0}},\boldsymbol{U}}[\mathbb{E}_{\boldsymbol{Z}}[\phi^{\boldsymbol{U}}_{j}(\boldsymbol{x_{0}}+\sigma\boldsymbol{Z},\boldsymbol{U})\mid\boldsymbol{x_{0}},\boldsymbol{U}]]\\
 & =\mathbb{E}_{\boldsymbol{x_{0}}\sim\boldsymbol{p_{0}},\boldsymbol{U}}[\mathbb{E}_{\xi\sim\mathcal{N}(0,1)}[\sum_{n=0}c_{n}(M^{\boldsymbol{U}}_{j}(\boldsymbol{x_{0}}),s_{j})\text{He}_{n}(\xi)]]\\
 & =\mathbb{E}_{\boldsymbol{x_{0}}\sim\boldsymbol{p_{0}},\boldsymbol{U}}[\sum_{n=0}c_{n}(M^{\boldsymbol{U}}_{j}(\boldsymbol{x_{0}}),s_{j})\mathbb{E}_{\xi\sim\mathcal{N}(0,1)}[\text{He}_{n}(\xi)\text{He}_{0}(\xi)]]\quad\quad(\text{since }\text{He}_{0}=1)\\
 & =\mathbb{E}_{\boldsymbol{x_{0}}\sim\boldsymbol{p_{0}},\boldsymbol{U}}[c_{0}(M^{\boldsymbol{U}}_{j}(\boldsymbol{x_{0}}),s_{j})]
\end{align*}

Rewrite $c_{0}(M^{\boldsymbol{U}}_{j}(\boldsymbol{x_{0}}),s_{j})\coloneqq g^{\boldsymbol{U}}_{j}(\boldsymbol{x_{0}})$
since $g^{\boldsymbol{U}}_{j}$ varies with $\boldsymbol{x_{0}}$
and $\boldsymbol{U}$ was given:
\begin{align*}
g^{\boldsymbol{U}}_{j}(\boldsymbol{x_{0}}) & =\mathbb{E}_{\xi\sim\mathcal{N}(0,1)}[f(\xi)\text{He}_{0}(\xi)\mid\boldsymbol{x_{0}},\boldsymbol{U}]\\
 & =\mathbb{E}_{\xi\sim\mathcal{N}(0,1)}[\omega(M^{\boldsymbol{U}}_{j}(\boldsymbol{x_{0}})+s_{j}\xi)]\\
 & =\int\omega(M^{\boldsymbol{U}}_{j}(\boldsymbol{x_{0}})+s_{j}\xi)\frac{e^{-\xi^{2}/2}}{\sqrt{2\pi}}d\xi
\end{align*}

Let $u=s_{j}\xi$, so $\xi=\frac{u}{s_{j}}$ and $d\xi=\frac{du}{s_{j}}$:
\begin{align*}
g^{\boldsymbol{U}}_{j}(\boldsymbol{x_{0}}) & =\int\omega(M^{\boldsymbol{U}}_{j}(\boldsymbol{x_{0}})+u)\frac{e^{-u^{2}/2s^{2}_{j}}}{s_{j}\sqrt{2\pi}}du\\
\end{align*}

This integral is the convolution of function $\omega$ and the Gaussian
$\mathcal{N}(0,s^{2}_{j})$, evaluated at the local point $M^{\boldsymbol{U}}_{j}(\boldsymbol{x_{0}})$,
i.e.
\begin{align*}
g^{\boldsymbol{U}}_{j}(\boldsymbol{x_{0}}) & =(\omega*\mathcal{N}_{s^{2}_{j}})(M^{\boldsymbol{U}}_{j}(\boldsymbol{x_{0}}))
\end{align*}

And so it follows that: 
\begin{equation}
\boldsymbol{\Sigma}^{\boldsymbol{U}}_{\phi,ij}=\text{Cov}_{\boldsymbol{x_{0}},\boldsymbol{U}}(g^{\boldsymbol{U}}_{i}(\boldsymbol{x_{0}}),g^{\boldsymbol{U}}_{j}(\boldsymbol{x_{0}}))+\mathbb{E}_{\boldsymbol{x_{0}},\boldsymbol{U}}[\sum_{n\geq1}n!\rho^{n}_{ij}c_{n}(M^{\boldsymbol{U}}_{i}(\boldsymbol{x_{0}}),s_{i})c_{n}(M^{\boldsymbol{U}}_{j}(\boldsymbol{x_{0}}),s_{j})]\label{eq:gen-var-phi}
\end{equation}

Now consider $\text{Cov}(\boldsymbol{x_{0}},\phi^{\boldsymbol{U}})_{ij}$
at row $i^{\text{th}}$ and column $j^{\text{th}}$:
\begin{align*}
\text{Cov}(\boldsymbol{x_{0}},\phi^{\boldsymbol{U}})_{ij} & =\mathbb{E}_{\boldsymbol{x_{0}},\boldsymbol{Z},\boldsymbol{U}}[(\boldsymbol{x}_{\boldsymbol{0},i}-\mu_{\boldsymbol{p_{0}},i})(\phi^{\boldsymbol{U}}(\boldsymbol{x_{0}}+\sigma\boldsymbol{Z}))_{j}]\\
 & =\mathbb{E}_{\boldsymbol{x_{0}},\boldsymbol{U}}[\mathbb{E}_{\boldsymbol{Z}\sim\mathcal{N}(0,\boldsymbol{I})}[(\boldsymbol{x}_{\boldsymbol{0},i}-\mu_{\boldsymbol{p_{0}},i})(\phi^{\boldsymbol{U}}(\boldsymbol{x_{0}}+\sigma\boldsymbol{Z},\boldsymbol{U}))_{j}\mid\boldsymbol{x_{0}},\boldsymbol{U}]]\\
 & =\mathbb{E}_{\boldsymbol{x_{0}},\boldsymbol{U}}[(\boldsymbol{x}_{\boldsymbol{0},i}-\mu_{\boldsymbol{p_{0}},i})\mathbb{E}_{\boldsymbol{Z}\sim\mathcal{N}(0,\boldsymbol{I})}[(\phi^{\boldsymbol{U}}(\boldsymbol{x_{0}}+\sigma\boldsymbol{Z},\boldsymbol{U}))_{j}\mid\boldsymbol{x_{0}},\boldsymbol{U}]]\\
 & =\mathbb{E}_{\boldsymbol{x_{0}},\boldsymbol{U}}[(\boldsymbol{x}_{\boldsymbol{0},i}-\mu_{\boldsymbol{p_{0}},i})g^{\boldsymbol{U}}_{j}(\boldsymbol{x_{0}})]
\end{align*}

Note that $g^{\boldsymbol{U}}_{j}(\boldsymbol{x_{0}})=c_{0}(M^{\boldsymbol{U}}_{j}(\boldsymbol{x_{0}}),s_{j})$
where $M^{\boldsymbol{U}}_{j}(\boldsymbol{x_{0}})=\theta^{\top}_{j}\boldsymbol{x_{0}}+\gamma^{\top}_{j}\boldsymbol{U}+\epsilon_{j}$
so really $g^{\boldsymbol{U}}_{j}(\boldsymbol{x_{0}})$ only depends
on $h_{j}\coloneqq\theta^{\top}_{j}\boldsymbol{x_{0}}+\gamma^{\top}_{j}\boldsymbol{U}+\epsilon_{j}$,
so rewrite $g^{\boldsymbol{U}}_{j}(\boldsymbol{x_{0}})=\psi(\theta^{\top}_{j}\boldsymbol{x_{0}}+\gamma^{\top}_{j}\boldsymbol{U}+\epsilon_{j})=\psi(h_{j})$
where $\psi(h)=c_{0}(h,s_{j})$. Now by Tower property:
\begin{align*}
\text{Cov}(\boldsymbol{x_{0}},\phi^{\boldsymbol{U}})_{ij} & =\mathbb{E}_{\boldsymbol{x_{0}},\boldsymbol{U}}[(\boldsymbol{x}_{\boldsymbol{0},i}-\mu_{\boldsymbol{p_{0}},i})g^{\boldsymbol{U}}_{j}(\boldsymbol{x_{0}})]\\
 & =\mathbb{E}_{h_{j}}[\mathbb{E}_{\boldsymbol{x_{0}},\boldsymbol{U}}[(\boldsymbol{x}_{\boldsymbol{0},i}-\mu_{\boldsymbol{p_{0}},i})g^{\boldsymbol{U}}_{j}(\boldsymbol{x_{0}})\mid h_{j}]]\\
 & =\mathbb{E}_{h_{j}}[g^{\boldsymbol{U}}_{j}(h_{j})\mathbb{E}_{\boldsymbol{x_{0}},\boldsymbol{U}}[(\boldsymbol{x}_{\boldsymbol{0},i}-\mu_{\boldsymbol{p_{0}},i})\mid h_{j}]]
\end{align*}

Define $r_{i}(h_{j})=\mathbb{E}_{\boldsymbol{x_{0}},\boldsymbol{U}}[(\boldsymbol{x}_{\boldsymbol{0},i}-\mu_{\boldsymbol{p_{0}},i})\mid h_{j}]$,
which is really a regression coefficient of $[\boldsymbol{x_{0}},\boldsymbol{U}]$
onto its own linear projection $h_{j}=[\theta_{j},\gamma_{j}]^{\top}[\boldsymbol{x_{0}},\boldsymbol{U}]$
plus noise. All this do is turning a multi-dimensional integral $\mathbb{E}_{\boldsymbol{x_{0}},\boldsymbol{U}}$
into an 1-D integral $\mathbb{E}_{h_{j}}$ .

$\mathcal{L}_{\sigma,\boldsymbol{U}}$ capture information of how
$\boldsymbol{x_{0}}$ relates to $\boldsymbol{U}$ in a non-linear
way.

Note: Recall that what follows is only valid when $\omega$ is component-wise,
so the formulas we derived would not apply to $\omega$ being a function
like softmax. In such case, $\mu^{\boldsymbol{U}}_{\phi}=\mathbb{E}[\phi^{\boldsymbol{U}}(\boldsymbol{x_{0}}+\sigma\boldsymbol{Z},\boldsymbol{U})]$,
$\boldsymbol{\Sigma}^{\boldsymbol{U}}_{\phi}=\text{Cov}(\phi(\boldsymbol{x_{0}}+\sigma\boldsymbol{Z},\boldsymbol{U}))$
and $\text{Cov}(\boldsymbol{x_{0}},\phi^{\boldsymbol{U}})=\mathbb{E}[(\boldsymbol{x_{0}}-\mu_{\boldsymbol{p_{0}}})(\phi^{\boldsymbol{U}}(\boldsymbol{x_{0}}+\sigma\boldsymbol{Z},\boldsymbol{U})-\mu^{\boldsymbol{U}}_{\phi})^{\top}]$
would have to be left in these forms and can't be broken down any
further, and remain integrals over all of $\boldsymbol{x_{0}}\sim\boldsymbol{p_{0}},\boldsymbol{Z},\boldsymbol{U}$.
Though in the case of softmax, if somehow the denominator can be asymptotically
approximated as some quantity that is independent of each of the entries
that sum to it, e.g. by LLN, then $\omega$ can still be considered
component-wise and our derivations apply.
\begin{align*}
I(X;U) & =\frac{1}{2}\int^{\infty}_{0}\frac{dt}{t^{2}}(\text{MMSE }[X\mid Y_{t}]-\text{MMSE }[X\mid Y_{t},U])\\
 & \approx\frac{1}{2}\int^{\infty}_{0}\frac{d\sigma}{\sigma^{3}}(\mathcal{L}_{\sigma}-\mathcal{L}_{\sigma,\boldsymbol{U}})\\
 & \approx\frac{1}{2}\int^{\infty}_{0}\frac{d\sigma}{\sigma^{3}}\text{Tr}(\text{Cov}(\boldsymbol{x_{0}},\phi^{\boldsymbol{U}})(\boldsymbol{\Sigma}^{\boldsymbol{U}}_{\phi})^{-1}\text{Cov}(\phi^{\boldsymbol{U}},\boldsymbol{x_{0}})-\text{Cov}(\boldsymbol{x_{0}},\phi)\boldsymbol{\Sigma}^{-1}_{\phi}\text{Cov}(\phi,\boldsymbol{x_{0}}))
\end{align*}

Note: $\text{Cov}(\boldsymbol{x_{0}},\phi)\boldsymbol{\Sigma}^{-1}_{\phi}\text{Cov}(\phi,\boldsymbol{x_{0}})$
is the explained variance of the best possible linear regression of
$\boldsymbol{x_{0}}$ onto $\phi$. 

so $\text{Cov}(\boldsymbol{x_{0}},\phi^{\boldsymbol{U}})(\boldsymbol{\Sigma}^{\boldsymbol{U}}_{\phi})^{-1}\text{Cov}(\phi^{\boldsymbol{U}},\boldsymbol{x_{0}})\geq\text{Cov}(\boldsymbol{x_{0}},\phi)\boldsymbol{\Sigma}^{-1}_{\phi}\text{Cov}(\phi,\boldsymbol{x_{0}})$
if $\text{span}(\phi)\subseteq\text{span}(\phi^{\boldsymbol{U}})$.
Otherwise, the integrand is not nonnegative.

\section{Jointly Gaussian $(\boldsymbol{x_{0}},\boldsymbol{U})$}

In the case of $(\boldsymbol{x_{0}},\boldsymbol{U})$ jointly Gaussian,
we can simplify the expression even more as this buys us Stein's lemma.
Now $l_{j}\coloneqq\theta^{\top}_{j}\boldsymbol{x_{0}}+\gamma^{\top}_{j}\boldsymbol{U}+\epsilon_{j}+s_{j}\xi$
is Gaussian, and equivalently $l_{j}=[\theta_{j},\gamma_{j}]^{\top}[\boldsymbol{x_{0}},\boldsymbol{U}]+\epsilon_{j}+s_{j}\xi$.
Let the co-variance between $\boldsymbol{x_{0}}$ and $\boldsymbol{U}$
be $C_{\boldsymbol{xU}}$. Then $l_{j}$ has mean and variance:
\begin{align*}
\tilde{m_{j}} & =\theta_{j}\mu_{\boldsymbol{x_{0}}}+\gamma_{j}\mu_{\boldsymbol{U}}\\
\tilde{S^{2}_{j}} & =\text{Var}([\theta_{j},\gamma_{j}]^{\top}[\boldsymbol{x_{0}},\boldsymbol{U}])+\text{Var}(s_{j}\xi)\\
 & =[\theta_{j},\gamma_{j}]^{\top}\begin{bmatrix}\boldsymbol{\Sigma}_{\boldsymbol{p_{0}}} & C_{\boldsymbol{xU}}\\
C_{\boldsymbol{Ux}} & \boldsymbol{\Sigma}_{\boldsymbol{U}}
\end{bmatrix}[\theta_{j},\gamma_{j}]+\sigma^{2}\|\theta_{j}\|\\
 & =\theta^{\top}_{j}\boldsymbol{\Sigma}_{\boldsymbol{p_{0}}}\theta_{j}+2\theta^{\top}_{j}C_{\boldsymbol{xU}}\gamma_{j}+\gamma^{\top}_{j}\boldsymbol{\Sigma}_{\boldsymbol{U}}\gamma_{j}+\sigma^{2}\|\theta_{j}\|
\end{align*}

Standardizing $l_{j}$ we get $l_{j}=\tilde{m_{j}}+\tilde{S_{j}}t$
where $t\sim\mathcal{N}(0,1)$. Now by Stein's lemma applying to $(\boldsymbol{x_{0}},l_{j})$
jointly Gaussian:
\begin{align*}
\text{Cov}(\boldsymbol{x_{0}},\phi^{\boldsymbol{U}})_{:j} & =\mathbb{E}_{\boldsymbol{x_{0}},\boldsymbol{Z},\boldsymbol{U}}[(\boldsymbol{x}_{\boldsymbol{0}}-\mu_{\boldsymbol{p_{0}}})(\omega(l_{j})]\\
 & =\text{Cov}(\boldsymbol{x_{0}},l_{j})\mathbb{E}[\omega'(l_{j})]\\
 & =(\boldsymbol{\Sigma}_{\boldsymbol{p_{0}}}\theta_{j}+C_{\boldsymbol{xU}}\gamma_{j})\mathbb{E}[\omega'(l_{j})]
\end{align*}

Rewrite $\omega(l_{j})=\omega(\tilde{m_{j}}+\tilde{S_{j}}t)\eqqcolon g(t)$,
so by Stein's lemma again:
\begin{align*}
\mathbb{E}_{t\sim\mathcal{N}(0,1)}[g'(t)] & =\mathbb{E}_{t\sim\mathcal{N}(0,1)}[t\cdot g(t)]
\end{align*}

So:
\begin{align*}
\mathbb{E}_{t\sim\mathcal{N}(0,1)}[\tilde{S_{j}}\omega'(\tilde{m_{j}}+\tilde{S_{j}}t)] & =\mathbb{E}_{t\sim\mathcal{N}(0,1)}[\text{He}_{1}\cdot\omega(\tilde{m_{j}}+\tilde{S_{j}}t)]\\
\tilde{S_{j}}\mathbb{E}_{l_{j}\sim\mathcal{N}(\tilde{m_{j}},\tilde{S^{2}_{j}})}[\omega'(l_{j})] & =c_{1}(\tilde{m_{j}},\tilde{S_{j}})
\end{align*}

Therefore:
\[
\mathbb{E}_{l_{j}\sim\mathcal{N}(\tilde{m_{j}},\tilde{S^{2}_{j}})}[\omega'(l_{j})]=\frac{c_{1}(\tilde{m_{j}},\tilde{S_{j}})}{\tilde{S_{j}}}
\]

and:
\[
\text{Cov}(\boldsymbol{x_{0}},\phi^{\boldsymbol{U}})_{:j}=(\boldsymbol{\Sigma}_{\boldsymbol{p_{0}}}\theta_{j}+C_{\boldsymbol{xU}}\gamma_{j})\alpha_{j}\hspace*{3em}\text{where }\alpha_{j}=\frac{c_{1}(\tilde{m_{j}},\tilde{S_{j}})}{\tilde{S_{j}}}
\]

Meanwhile, we consider $\boldsymbol{\Sigma}^{\boldsymbol{U}}_{\phi,ij}$;
first standardizing each of $l_{i}$ and $l_{j}$ is:
\[
t_{i}=\frac{l_{i}-\tilde{m_{i}}}{\tilde{S_{i}}}\sim\mathcal{N}(0,1)\hspace*{5em}t_{j}=\frac{l_{j}-\tilde{m_{j}}}{\tilde{S_{j}}}\sim\mathcal{N}(0,1)
\]

and their correlation is:
\[
\tilde{r}_{ij}=\text{Corr}(t_{i},t_{j})=\frac{\theta^{\top}_{i}\boldsymbol{\Sigma}_{\boldsymbol{p_{0}}}\theta_{j}+\theta^{\top}_{i}C_{\boldsymbol{xU}}\gamma_{j}+\gamma^{\top}_{i}C_{\boldsymbol{Ux}}\theta_{j}+\gamma^{\top}_{i}\boldsymbol{\Sigma}_{\boldsymbol{U}}\gamma_{j}+\sigma^{2}\theta^{\top}_{i}\theta_{j}}{\tilde{S_{i}}\tilde{S_{j}}}
\]

Hermite-expansion:
\begin{align*}
\omega(l_{i})=\omega(\tilde{m_{i}}+\tilde{S_{i}}t_{i}) & =\sum_{n=0}c_{n}(\tilde{m_{i}},\tilde{S_{i}})\text{He}_{n}(t_{i})\\
\omega(l_{j})=\omega(\tilde{m_{j}}+\tilde{S_{j}}t_{j}) & =\sum_{m=0}c_{m}(\tilde{m_{j}},\tilde{S_{j}})\text{He}_{m}(t_{j})
\end{align*}

Then:
\begin{align*}
\mathbb{E}[\phi^{\boldsymbol{U}}_{i}\phi^{\boldsymbol{U}}_{j}] & =\sum_{n,m}c_{n}(\tilde{m_{i}},\tilde{S_{i}})\text{He}_{n}(t_{i})c_{m}(\tilde{m_{j}},\tilde{S_{j}})\text{He}_{m}(t_{j})\mathbb{E}\big[\text{He}_{n}(t_{i})\text{He}_{m}(t_{j})\big]\\
 & =\sum_{n}n!\tilde{r}^{n}_{ij}c_{n}(\tilde{m_{i}},\tilde{S_{i}})c_{n}(\tilde{m_{j}},\tilde{S_{j}})
\end{align*}

And so:
\begin{align*}
\boldsymbol{\Sigma}^{\boldsymbol{U}}_{\phi,ij} & =\mathbb{E}[\phi^{\boldsymbol{U}}_{i}\phi^{\boldsymbol{U}}_{j}]-\mathbb{E}[\phi^{\boldsymbol{U}}_{i}]\mathbb{E}[\phi^{\boldsymbol{U}}_{j}]\\
 & =\sum_{n}n!\tilde{r}^{n}_{ij}c_{n}(\tilde{m_{i}},\tilde{S_{i}})c_{n}(\tilde{m_{j}},\tilde{S_{j}})-c_{0}(\tilde{m_{i}},\tilde{S_{i}})c_{0}(\tilde{m_{j}},\tilde{S_{j}})\\
 & =\sum_{n\geq1}n!\tilde{r}^{n}_{ij}c_{n}(\tilde{m_{i}},\tilde{S_{i}})c_{n}(\tilde{m_{j}},\tilde{S_{j}})
\end{align*}

This only has 1 term instead of 2 like in the general case, since
$l_{i}$ and $l_{j}$ are both Gaussians, so we can write a closed-form
for their covariance instead of having to decompose it out like in
(\ref{eq:gen-var-phi}).

So writing $\boldsymbol{A}_{\boldsymbol{U}}=(\boldsymbol{\Sigma}_{\boldsymbol{p_{0}}}\Theta^{\top}+C_{\boldsymbol{xU}}\Gamma^{\top})\text{diag}(\alpha)$,
then:
\[
\mathcal{L}_{\sigma,\boldsymbol{U}}=\text{Tr}(\boldsymbol{\Sigma}_{\boldsymbol{p_{0}}})-\text{Tr}(\boldsymbol{A}_{\boldsymbol{U}}(\boldsymbol{\Sigma}^{\boldsymbol{U}}_{\phi})^{-1}\boldsymbol{A}^{\top}_{\boldsymbol{U}})
\]

Similarly for the unconditional case:
\begin{align*}
m_{j} & =\theta_{j}\mu_{\boldsymbol{x_{0}}}\\
S^{2}_{j} & =\theta^{\top}_{j}\boldsymbol{\Sigma}_{\boldsymbol{p_{0}}}\theta_{j}+\sigma^{2}\|\theta_{j}\|
\end{align*}

and:
\[
\text{Cov}(\boldsymbol{x_{0}},\phi)_{:j}=(\boldsymbol{\Sigma}_{\boldsymbol{p_{0}}}\theta_{j})\alpha^{0}_{j}\hspace*{3em}\text{where }\alpha^{0}_{j}=\frac{c_{1}(m_{j},S_{j})}{S_{j}}
\]
\[
\boldsymbol{\Sigma}_{\phi,ij}=\sum_{n\geq1}n!r^{n}_{ij}c_{n}(m_{i},S_{i})c_{n}(m_{j},S_{j})
\]

And so:
\begin{align*}
\mathcal{L}_{\sigma}-\mathcal{L}_{\sigma,\boldsymbol{U}} & =\text{Tr}(\boldsymbol{A}_{\boldsymbol{U}}(\boldsymbol{\Sigma}^{\boldsymbol{U}}_{\phi})^{-1}\boldsymbol{A}^{\top}_{\boldsymbol{U}})-\text{Tr}(\boldsymbol{A}(\boldsymbol{\Sigma}_{\phi})^{-1}\boldsymbol{A}^{\top})
\end{align*}

where $\boldsymbol{A}=\boldsymbol{\Sigma}_{\boldsymbol{p_{0}}}\Theta^{\top}\text{diag}(\alpha^{0})$ 
\end{document}
